\FloatBarrier
\section{Esercitazioni di Clustering Eterogeneo e Multi-Vista}
\label{ch:dm-ex-heterogeneous-multiview}

\subsection{Convenzioni Operative}
\label{sec:dm-ex-hmv-rules}

Negli esercizi di \textit{clustering} su dati categorici, misti e multi-vista valgono le seguenti convenzioni:
\begin{itemize}
    \item \textbf{Spareggio per le mode categoriche:} in caso di parit\`a di frequenza, si sceglie la categoria che precede in \textbf{ordine alfabetico}.
    \item \textbf{Spareggio nell'assegnazione ai cluster:} in caso di parit\`a di dissimilarit\`a rispetto a due rappresentanti, si assegna l'oggetto al cluster con \textbf{indice minore}.
    \item \textbf{Assegnazioni e metriche:} tutte le assegnazioni sono \textit{hard}; si usa la distanza Euclidea al quadrato per gli attributi numerici e la dissimilarit\`a di corrispondenza semplice (\textit{simple matching}) per gli attributi categorici.
\end{itemize}

% ====================================================================
\FloatBarrier
\subsection{Esercizio 1: K-Modes}
\label{sec:dm-ex-kmodes}

Si considerino $6$ oggetti ($A, \dots, F$) descritti da $m = 3$ attributi categorici (\texttt{Color}, \texttt{Shape}, \texttt{Size}), con $k = 2$ cluster e mode iniziali (\autoref{tab:dm-ex-kmodes}):
\begin{equation*}
    Q_1^{(0)} = (\text{Red}, \text{Circle}, \text{Small}), \qquad Q_2^{(0)} = (\text{Blue}, \text{Square}, \text{Large})
\end{equation*}

\begin{table}[H]
\centering
\caption{Dataset e prima assegnazione nell'esercizio su K-Modes.}
\label{tab:dm-ex-kmodes}
\small
\begin{tabular}{ccccccc}
\toprule
\textbf{Object} & \textbf{Color} & \textbf{Shape} & \textbf{Size} & $d(X_i, Q_1^{(0)})$ & $d(X_i, Q_2^{(0)})$ & \textbf{Cluster} \\
\midrule
$A$ & Red  & Circle & Small & $0$ & $3$ & $1$ \\
$B$ & Red  & Square & Small & $1$ & $2$ & $1$ \\
$C$ & Blue & Circle & Large & $2$ & $1$ & $2$ \\
$D$ & Blue & Square & Large & $3$ & $0$ & $2$ \\
$E$ & Red  & Circle & Large & $1$ & $2$ & $1$ \\
$F$ & Blue & Circle & Small & $1$ & $2$ & $1$ \\
\bottomrule
\end{tabular}
\end{table}

\begin{enumerate}
    \item \textbf{Dissimilarit\`a e assegnazione:} usando $d(X_i, Q_l) = \sum_{j=1}^m \mathbf{1}(x_{ij} \neq q_{lj})$, si ottengono i cluster $C_1 = \{A, B, E, F\}$ e $C_2 = \{C, D\}$.
    \item \textbf{Aggiornamento delle mode:} in $C_1$, le frequenze massime sono \texttt{Red} ($3/4$), \texttt{Circle} ($3/4$) e \texttt{Small} ($3/4$), quindi $Q_1^{(1)} = (\text{Red}, \text{Circle}, \text{Small})$. In $C_2$, \texttt{Blue} ($2/2$) e \texttt{Large} ($2/2$) sono unanimi, mentre su \texttt{Shape} c'\`e parit\`a tra \texttt{Circle} e \texttt{Square}; lo spareggio alfabetico seleziona \texttt{Circle}, quindi $Q_2^{(1)} = (\text{Blue}, \text{Circle}, \text{Large})$.
    \item \textbf{Funzione obiettivo e stabilit\`a:}
    \begin{equation*}
        P(W, Q^{(1)}) = \sum_{l=1}^2 \sum_{X_i \in C_l} d(X_i, Q_l^{(1)}) = (0 + 1 + 1 + 1) + (0 + 1) = 4
    \end{equation*}
    Al passo successivo, per $F$ si ha $d(F, Q_1^{(1)}) = 1 = d(F, Q_2^{(1)})$; per la regola dell'indice minore $F$ resta in $C_1$ e tutte le assegnazioni restano invariate: il clustering \`e \textbf{stabile}.
\end{enumerate}

% ====================================================================
\FloatBarrier
\subsection{Esercizio 2: K-Prototypes}
\label{sec:dm-ex-kprototypes}

Si considerino $6$ oggetti descritti da due attributi numerici (\texttt{Age}, \texttt{Income}) e uno categorico (\texttt{Job}), con $k = 2$, peso $\gamma = 8$ e prototipi iniziali $Q_1^{(0)} = (22, 22, \text{Student})$ e $Q_2^{(0)} = (32, 32, \text{Engineer})$ (\autoref{tab:dm-ex-kproto}).

\begin{table}[H]
\centering
\caption{Dataset e assegnazione iniziale nell'esercizio su K-Prototypes ($\gamma = 8$).}
\label{tab:dm-ex-kproto}
\small
\begin{tabular}{ccccccc}
\toprule
\textbf{Object} & \textbf{Age} & \textbf{Income} & \textbf{Job} & $d(X_i, Q_1^{(0)})$ & $d(X_i, Q_2^{(0)})$ & \textbf{Cluster} \\
\midrule
$A$ & $20$ & $20$ & Student  & $8$   & $296$ & $1$ \\
$B$ & $22$ & $22$ & Student  & $0$   & $208$ & $1$ \\
$C$ & $24$ & $24$ & Engineer & $16$  & $128$ & $1$ \\
$D$ & $30$ & $30$ & Engineer & $136$ & $8$   & $2$ \\
$E$ & $32$ & $32$ & Engineer & $208$ & $0$   & $2$ \\
$F$ & $34$ & $34$ & Teacher  & $296$ & $16$  & $2$ \\
\bottomrule
\end{tabular}
\end{table}

\begin{enumerate}
    \item \textbf{Dissimilarit\`a mista:}
    \begin{equation*}
        d(X_i, Q_l) = (\text{Age}_i - \mu_{l,\text{Age}})^2 + (\text{Income}_i - \mu_{l,\text{Income}})^2 + \gamma\,\mathbf{1}(\text{Job}_i \neq \text{Job}_l)
    \end{equation*}
    Ad esempio, $d(C, Q_1^{(0)}) = 2^2 + 2^2 + 8 = 16$.
    \item \textbf{Assegnazione e aggiornamento:} si ottengono $C_1 = \{A,B,C\}$ e $C_2 = \{D,E,F\}$. Aggiornando le coordinate numeriche con la media e quella categorica con la moda, si ha $Q_1^{(1)} = (22, 22, \text{Student})$ e $Q_2^{(1)} = (32, 32, \text{Engineer})$ (convergenza raggiunta).
    \item \textbf{Funzione obiettivo:} $P(W, Q^{(1)}) = (8 + 0 + 16) + (8 + 0 + 16) = 48$.
\end{enumerate}

% ====================================================================
\FloatBarrier
\subsection{Esercizio 3: Early Fusion --- Distanze e Pesi delle Viste}
\label{sec:dm-ex-early-fusion}

Si considerino $4$ oggetti descritti da due viste numeriche bidimensionali:
\begin{itemize}
    \item \textbf{Vista 1:} $x_A^{(1)} = (0,0)$, $x_B^{(1)} = (1,1)$, $x_C^{(1)} = (3,0)$, $x_D^{(1)} = (2,3)$;
    \item \textbf{Vista 2:} $x_A^{(2)} = (0,0)$, $x_B^{(2)} = (3,4)$, $x_C^{(2)} = (1,0)$, $x_D^{(2)} = (4,4)$.
\end{itemize}

\begin{enumerate}
    \item \textbf{Rappresentazione \textit{Early Fusion} ($x_i = [x_i^{(1)} \mid x_i^{(2)}]$):} $x_A = (0,0,0,0)$, $x_B = (1,1,3,4)$, $x_C = (3,0,1,0)$.
    \item \textbf{Distanze Euclidee al quadrato non pesate:} $d^2(A,B) = (1^2+1^2) + (3^2+4^2) = 2 + 25 = 27$; $d^2(A,C) = (3^2+0^2) + (1^2+0^2) = 9 + 1 = 10$. L'oggetto pi\`u vicino ad $A$ \`e \textbf{$C$}.
    \item \textbf{Distanze pesate per vista con $(\alpha_1, \alpha_2) = (4, 0.2)$:}
    \begin{equation*}
        d_\alpha^2(A, B) = 4(2) + 0.2(25) = 13, \qquad d_\alpha^2(A, C) = 4(9) + 0.2(1) = 36.2
    \end{equation*}
    L'oggetto pi\`u vicino ad $A$ diventa \textbf{$B$}: la pesatura delle viste pu\`o modificare radicalmente la geometria dello spazio e l'esito del clustering.
    \item \textbf{Standardizzazione vs.\ pesatura delle viste:} la standardizzazione riscala separatamente ciascuna singola feature a varianza unitaria, ma poich\'e la distanza Euclidea somma un termine quadratico per ogni feature, una vista con molte pi\`u feature continua a dominare la distanza complessiva se non viene introdotta una pesatura o normalizzazione a livello di vista.
\end{enumerate}

% ====================================================================
\FloatBarrier
\subsection{Esercizio 4: MCLA --- Ipergrafo e Meta-Grafo}
\label{sec:dm-ex-mcla-graph}

Su $n = 6$ oggetti $\{A, B, C, D, E, F\}$, due partizioni producono gli iperarchi $e_1 = \{A,B,C\}$, $e_2 = \{D,E,F\}$ (Partizione 1) ed $e_3 = \{A,B\}$, $e_4 = \{C,D\}$, $e_5 = \{E,F\}$ (Partizione 2).

\begin{enumerate}
    \item \textbf{Matrice di incidenza $H \in \{0,1\}^{6 \times 5}$} (righe $A \dots F$, colonne $e_1 \dots e_5$) e \textbf{matrice di adiacenza pesata $W$} del meta-grafo $G_M$ ($W_{ab} = J(e_a, e_b)$ per $a \neq b$):
    \begin{equation*}
        H = \begin{pmatrix}
        1 & 0 & 1 & 0 & 0 \\
        1 & 0 & 1 & 0 & 0 \\
        1 & 0 & 0 & 1 & 0 \\
        0 & 1 & 0 & 1 & 0 \\
        0 & 1 & 0 & 0 & 1 \\
        0 & 1 & 0 & 0 & 1
        \end{pmatrix}, \qquad
        W = \begin{pmatrix}
        0 & 0 & \frac{2}{3} & \frac{1}{4} & 0 \\
        0 & 0 & 0 & \frac{1}{4} & \frac{2}{3} \\
        \frac{2}{3} & 0 & 0 & 0 & 0 \\
        \frac{1}{4} & \frac{1}{4} & 0 & 0 & 0 \\
        0 & \frac{2}{3} & 0 & 0 & 0
        \end{pmatrix}
    \end{equation*}
    \item \textbf{Similarit\`a di Jaccard:} tramite $J(e_a, e_b) = \frac{h_a^\top h_b}{\norm{h_a}_2^2 + \norm{h_b}_2^2 - h_a^\top h_b}$ si ottengono $J(e_1, e_3) = \frac{2}{3}$, $J(e_1, e_4) = \frac{1}{4}$, $J(e_2, e_4) = \frac{1}{4}$, $J(e_2, e_5) = \frac{2}{3}$.
    \item \textbf{Struttura bipartita:} all'interno della medesima partizione \textit{hard} i cluster sono disgiunti ($J = 0$), quindi archi di peso positivo possono collegare solo iperarchi di partizioni diverse: $G_M$ \`e \textbf{bipartito}.
\end{enumerate}

% ====================================================================
\FloatBarrier
\subsection{Esercizio 5: MCLA --- Collasso dei Meta-Cluster e Compete for Objects}
\label{sec:dm-ex-mcla-compete}

Dati i due meta-cluster $\mathcal{M}_1 = \{e_1, e_3\}$ e $\mathcal{M}_2 = \{e_2, e_4, e_5\}$, si calcolano i vettori di associazione medi $a_j = \frac{1}{\abs{\mathcal{M}_j}} \sum_{e_a \in \mathcal{M}_j} h_a$:
\begin{align*}
    a_1 &= \frac{h_1 + h_3}{2} = \left(1, 1, \frac{1}{2}, 0, 0, 0\right)^\top, \\
    a_2 &= \frac{h_2 + h_4 + h_5}{3} = \left(0, 0, \frac{1}{3}, \frac{2}{3}, \frac{2}{3}, \frac{2}{3}\right)^\top
\end{align*}

Applicando la regola \textit{Compete for Objects} ($\lambda^*(x_i) = \arg\max_j A_{ij}$ sulla matrice $A = [a_1 \mid a_2]$, \autoref{tab:dm-ex-mcla-assoc}), $A, B, C$ favoriscono $\mathcal{M}_1$ (per $C$, $\frac{1}{2} > \frac{1}{3}$) mentre $D, E, F$ favoriscono $\mathcal{M}_2$, generando la partizione di consenso $C_1^* = \{A,B,C\}$, $C_2^* = \{D,E,F\}$. Poich\'e ciascuna colonna media le frazioni di appartenenza dentro meta-cluster di cardinalit\`a diversa, le righe di $A$ non sommano necessariamente a $1$ (es.\ per $C$: $\frac{1}{2} + \frac{1}{3} = \frac{5}{6}$): i valori $A_{ij}$ rappresentano forze di associazione e non probabilit\`a normalizzate.

\begin{table}[H]
\centering
\caption{Matrice di associazione $A = [a_1 \mid a_2]$ e assegnazione finale di consenso.}
\label{tab:dm-ex-mcla-assoc}
\small
\begin{tabular}{cccc}
\toprule
\textbf{Object} & $\mathcal{M}_1$ & $\mathcal{M}_2$ & \textbf{Consenso} \\
\midrule
$A$ & $1$           & $0$           & $C_1^*$ \\
$B$ & $1$           & $0$           & $C_1^*$ \\
$C$ & $\frac{1}{2}$ & $\frac{1}{3}$ & $C_1^*$ \\
$D$ & $0$           & $\frac{2}{3}$ & $C_2^*$ \\
$E$ & $0$           & $\frac{2}{3}$ & $C_2^*$ \\
$F$ & $0$           & $\frac{2}{3}$ & $C_2^*$ \\
\bottomrule
\end{tabular}
\end{table}

% ====================================================================
\FloatBarrier
\subsection{Esercizio 6: Multi-View K-Means --- Aggiornamenti Alternati}
\label{sec:dm-ex-mv-kmeans}

Su $6$ oggetti descritti da due viste numeriche 1D (\autoref{tab:dm-ex-mv-data}), con partizione iniziale nella Vista 1 pari a $C_1^{(1)} = \{A,B,C\}$ e $C_2^{(1)} = \{D,E,F\}$, si esegue un aggiornamento alternato completo:

\begin{enumerate}
    \item \textbf{Centroidi nella Vista 2} dalle appartenenze della Vista 1:
    \begin{equation*}
        \mu_1^{(2)} = \frac{0 + 1 + 6}{3} = \frac{7}{3}, \qquad \mu_2^{(2)} = \frac{4 + 8 + 9}{3} = 7
    \end{equation*}
    \item \textbf{Riassegnazione nella Vista 2} (\autoref{tab:dm-ex-mv-data}): $C$ ($6$) si sposta nel cluster $2$ e $D$ ($4$) nel cluster $1$, formando $C_1^{(2)} = \{A,B,D\}$ e $C_2^{(2)} = \{C,E,F\}$.
    \item \textbf{Centroidi nella Vista 1} dalle appartenenze della Vista 2:
    \begin{equation*}
        \mu_1^{(1)} = \frac{0 + 1 + 7}{3} = \frac{8}{3}, \qquad \mu_2^{(1)} = \frac{2 + 8 + 9}{3} = \frac{19}{3}
    \end{equation*}
    \item \textbf{Riassegnazione nella Vista 1:} con punto medio $\frac{8/3 + 19/3}{2} = 4.5$, gli oggetti $A,B,C$ ($< 4.5$) tornano nel cluster $1$ e $D,E,F$ ($> 4.5$) nel cluster $2$: $C_1^{(1),\text{new}} = \{A,B,C\}$, $C_2^{(1),\text{new}} = \{D,E,F\}$.
\end{enumerate}

\begin{table}[H]
\centering
\caption{Dati Multi-View e riassegnazione degli oggetti nella Vista 2.}
\label{tab:dm-ex-mv-data}
\small
\resizebox{\columnwidth}{!}{%
\begin{tabular}{cccccc}
\toprule
\textbf{Object} & \textbf{View 1: $x_i^{(1)}$} & \textbf{View 2: $x_i^{(2)}$} & $d^2$ da $\mu_1^{(2)} = \frac{7}{3}$ & $d^2$ da $\mu_2^{(2)} = 7$ & \textbf{Cluster (Vista 2)} \\
\midrule
$A$ & $0$ & $0$ & $\frac{49}{9}$  & $49$ & $1$ \\
$B$ & $1$ & $1$ & $\frac{16}{9}$  & $36$ & $1$ \\
$C$ & $2$ & $6$ & $\frac{121}{9}$ & $1$  & $2$ \\
$D$ & $7$ & $4$ & $\frac{25}{9}$  & $9$  & $1$ \\
$E$ & $8$ & $8$ & $\frac{289}{9}$ & $1$  & $2$ \\
$F$ & $9$ & $9$ & $\frac{400}{9}$ & $4$  & $2$ \\
\bottomrule
\end{tabular}%
}
\end{table}

Il ritorno della Vista 1 alla partizione iniziale non rende ridondante il passo sulla Vista 2 (dove $C$ e $D$ hanno scambiato cluster): \textit{Multi-View K-Means} scambia tra le viste le \textbf{appartenenze oggetto-cluster} e non le coordinate dei centroidi (ciascuna vista calcola i propri centroidi nel proprio spazio), e un solo aggiornamento alternato non implica ancora la convergenza globale.
